% !TEX root = ../main.tex
\begin{ouverture}[piles]{14}{Énergie et puissance électriques}{Partie 3 · L'énergie : conversions et transferts}
{Intensité et charges\enspace·\enspace Source réelle\enspace·\enspace Puissance\enspace·\enspace Effet Joule\enspace·\enspace Rendement}
\objectif{Relier l'intensité d'un courant continu au débit de charges électriques.}
\objectif{Modéliser une source réelle de tension par une source idéale en série avec une résistance ; exploiter sa caractéristique.}
\objectif{Établir un bilan de puissance dans un circuit ; calculer la puissance dissipée par effet Joule.}
\objectif{Définir et évaluer le rendement d'un convertisseur ; citer des ordres de grandeur de puissances.}
\prerequis{Tension, intensité, loi d'Ohm $U=R\,I$, caractéristique d'un dipôle (classe de seconde).}
\prerequis{Énergie, puissance, relation $E=P\,\Delta t$, conversions d'énergie (collège).}
\end{ouverture}
\essentiel
\essSection{1}{Intensité et porteurs de charge}
\begin{definition}[Intensité d'un courant continu]
Un courant électrique est un déplacement d'ensemble de \textbf{porteurs de charge} : électrons libres dans les métaux, ions dans les
solutions. L'intensité est le \textbf{débit de charges} à travers une section du circuit :
\[I=\frac{Q}{\Delta t}\qquad\text{($I$ en A, $Q$ en C, $\Delta t$ en s).}\]
Un courant de \qty{0,50}{A} fait passer \qty{30}{C} en une minute, soit $\dfrac{30}{\num{1,60e-19}}\approx\num{1,9e20}$ électrons.
\end{definition}
\schema{14-porteurs}
\essSection{2}{Modèle d'une source réelle de tension}
\begin{loi}[Source réelle de tension continue]
Une source réelle (pile, batterie) se modélise par une source idéale de tension $E$ en série avec une résistance interne $r$. La
tension entre ses bornes diminue quand l'intensité débitée augmente :
\[U=E-r\,I.\]
La caractéristique $U=f(I)$ est une droite d'ordonnée à l'origine $E$ et de pente $-r$. En court-circuit ($U=0$), l'intensité vaut $I_\mathrm{cc}=\dfrac{E}{r}$.
Exemple : avec $E=\qty{4,5}{V}$ et $r=\qty{1,5}{\ohm}$, la tension tombe à \qty{3,9}{V} pour \qty{0,40}{A}.
\end{loi}
\par\penalty0
\schema{14-modele}
\begin{python}
import numpy as np
I = [0, 0.10, 0.20, 0.30, 0.40, 0.50]         # A
U = [4.52, 4.36, 4.21, 4.07, 3.90, 3.76]      # V : pile mesuree
a, b = np.polyfit(I, U, 1)                    # droite U = a I + b
print('E =', round(b, 2), 'V ; r =', round(-a, 2), 'ohm')
\end{python}
\begin{console}
E = 4.52 V ; r = 1.52 ohm
\end{console}
\noindent\capacite{Capacité numérique} Le programme modélise les mesures par une droite : l'ordonnée à l'origine donne $E$, l'opposé de la
pente donne $r$.
\essSection{3}{Puissance, énergie et effet Joule}
\begin{definition}[Puissance et énergie électriques]
Un dipôle soumis à la tension $U$ et traversé par l'intensité $I$ reçoit (ou fournit) la puissance $P=U\,I$ (en W). Pendant $\Delta t$, l'énergie
transférée vaut $E=P\,\Delta t$ (en J ; $\qty{1}{kWh}=\qty{3,6e6}{J}$).
\end{definition}
\begin{loi}[Effet Joule et bilan de puissance]
Un conducteur ohmique de résistance $R$ convertit toute l'énergie reçue en énergie thermique : $P_\mathrm{J}=R\,I^2=\dfrac{U^2}{R}$. Pour une source
réelle, $E\,I=U\,I+r\,I^2$ : la puissance $E\,I$ issue de la conversion (chimique, par exemple) se partage entre la puissance $U\,I$
fournie au circuit et la puissance $r\,I^2$ dissipée dans la source. Exemple : à \qty{0,40}{A}, $E\,I=\qty{1,80}{W}$, $U\,I=\qty{1,56}{W}$ et $r\,I^2=\qty{0,24}{W}$.
\end{loi}
\begin{definition}[Rendement d'un convertisseur]
Un convertisseur (moteur, panneau solaire, lampe…) reçoit une puissance $P_\mathrm{reçue}$ et en restitue une partie utile :
\[\eta=\frac{P_\mathrm{utile}}{P_\mathrm{reçue}}<1,\qquad\text{le reste étant dissipé, le plus souvent sous forme thermique.}\]
\end{definition}
\par\penalty0
\schema{14-bilan}
\begin{propriete}[Quelques ordres de grandeur de puissances]
\par\smallskip\noindent\hfill\begin{tabular}{lr@{\qquad}lr}
\toprule
Dispositif & Puissance & Dispositif & Puissance \\
\midrule
Lampe à LED & \qty{5}{W} à \qty{10}{W} & Bouilloire, four & \qty{2}{kW} à \qty{3}{kW} \\
Chargeur de smartphone & \qty{5}{W} à \qty{30}{W} & Voiture électrique & \qty{100}{kW} \\
Ordinateur portable & \qty{50}{W} & Rame de TGV & \qty{9}{MW} \\
Réfrigérateur (en marche) & \qty{150}{W} & Réacteur nucléaire & \qty{1}{GW} \\
\bottomrule
\end{tabular}\hfill\null\par\smallskip
\end{propriete}
\begin{attention}
Ne pas confondre $E$, la tension de la source idéale (en V), et l'énergie (en J). La tension aux bornes d'une pile n'est égale à $E$ que
si elle ne débite pas. En court-circuit, toute la puissance est dissipée dans la pile : elle chauffe dangereusement.
\end{attention}
\begin{experience}[Caractéristique d'une pile et rendement d'un moteur]
On fait varier l'intensité débitée par une pile grâce à une résistance réglable et l'on relève $U$ et $I$ : la droite obtenue donne $E$
et $r$. Un petit moteur soulève une masse connue : on compare la puissance mécanique $\dfrac{m\,g\,h}{\Delta t}$ à la puissance électrique $U\,I$.
\end{experience}
\begin{rediger}
Écrire la relation littérale, préciser le dipôle étudié, convertir en unités SI (mA en A, mAh en C, min en s, kW en W),
et ne jamais confondre puissance (W) et énergie (J ou kWh).
\end{rediger}
\exercices

\rubrique{Découverte}
\begin{blocLivre}{ex-14-1}
\begin{exercice}[1]{Un débit d'électrons}
Une lampe de poche est traversée par un courant d'intensité $I=\qty{0,20}{A}$ pendant une minute.
\begin{donnees}
$e=\qty{1,60e-19}{C}$.
\end{donnees}
\begin{questions}
\item Quels sont les porteurs de charge dans les fils métalliques ? Dans quel sens se déplacent-ils ?
\item Calculer la charge qui traverse une section du fil en une minute.
\item En déduire le nombre d'électrons correspondant.
\item Pourquoi la lampe s'allume-t-elle dès que l'on ferme l'interrupteur, alors que les électrons avancent très lentement ?
\end{questions}
\end{exercice}
\end{blocLivre}
\begin{blocLivre}{ex-14-2}
\begin{exercice}[1]{Lire une caractéristique}
\schema{14-ex02}
\begin{questions}
\item Déterminer graphiquement la valeur de $E$.
\item Déterminer la résistance interne $r$ de la pile et écrire la relation entre $U$ et $I$.
\item Quelle tension la pile délivre-t-elle quand elle débite \qty{0,40}{A} ? Quelle puissance fournit-elle alors ?
\item Que vaut l'intensité de court-circuit ? Pourquoi ne faut-il jamais court-circuiter une pile ?
\end{questions}
\end{exercice}
\end{blocLivre}
\begin{blocLivre}{ex-14-3}
\begin{exercice}[1]{Puissance et énergie}
Une bouilloire de \qty{2,0}{kW} chauffe l'eau d'un thé en \qty{3}{min}. Le kilowattheure coûte 0,25~€.
\begin{questions}
\item Calculer l'énergie consommée, en joules puis en kilowattheures.
\item Calculer le coût de l'opération.
\item Une lampe à LED de \qty{7}{W} éclaire autant qu'une ancienne lampe à incandescence de \qty{60}{W}. Calculer l'énergie consommée par
chacune en \qty{1000}{h}.
\item Quelle économie la lampe à LED permet-elle sur cette durée ?
\end{questions}
\end{exercice}
\end{blocLivre}
\begin{blocLivre}{ex-14-4}
\begin{exercice}[1]{Un conducteur ohmique qui chauffe}
Un conducteur ohmique de résistance $R=\qty{47}{\ohm}$ est traversé par un courant d'intensité \qty{0,10}{A}. Le fabricant indique qu'il ne
doit pas dissiper plus de \qty{0,25}{W}.
\begin{questions}
\item Calculer la tension à ses bornes.
\item Calculer la puissance dissipée par effet Joule. Sous quelle forme l'énergie est-elle transférée ?
\item Calculer l'énergie dissipée en \qty{10}{min}.
\item Le conducteur risque-t-il d'être détérioré ? Quelle intensité maximale peut-il supporter ?
\end{questions}
\end{exercice}
\end{blocLivre}
\begin{blocLivre}{ex-14-5}
\begin{exercice}[1]{Le rendement d'un moteur}
Un petit moteur alimenté sous \qty{12}{V} est traversé par un courant de \qty{2,5}{A}. Il soulève une charge et fournit une puissance
mécanique de \qty{21}{W}.
\begin{questions}
\item Calculer la puissance électrique reçue par le moteur.
\item Calculer son rendement.
\item Représenter la chaîne de puissance du moteur et préciser la forme des pertes.
\item Calculer l'énergie perdue en \qty{5}{min} de fonctionnement.
\end{questions}
\end{exercice}
\end{blocLivre}
\rubrique{Entraînement}
\begin{blocLivre}{ex-14-6}
\begin{exercice}[2]{La facture d'électricité, en Python}
\begin{python}
appareils = {'box internet': (12, 24), 'television': (90, 4),
             'refrigerateur': (150, 8), 'four': (2500, 0.5)}
total = 0
for nom, (P, h) in appareils.items():
    E = P * h * 365 / 1000            # kWh par an
    total = total + E
    print(nom, ':', round(E), 'kWh')
print('total :', round(total), 'kWh, soit', round(total * 0.25), 'euros')
\end{python}
\begin{console}
box internet : 105 kWh
television : 131 kWh
refrigerateur : 438 kWh
four : 456 kWh
total : 1131 kWh, soit 283 euros
\end{console}
\begin{questions}
\item \capacite{Capacité numérique}Que représentent les deux nombres associés à chaque appareil ? Quelle relation la ligne \code{E = ...} utilise-t-elle ?
\item Pourquoi compte-t-on seulement 8 heures par jour pour le réfrigérateur, qui reste toujours branché ?
\item La box, de faible puissance, coûte-t-elle peu ? Calculer son coût annuel.
\item Une box mise en veille consomme encore \qty{6}{W}. Calculer l'énergie consommée en veille, 20 heures par jour, pendant un an.
\end{questions}
\end{exercice}
\end{blocLivre}
\begin{blocLivre}{ex-14-7}
\begin{exercice}[2]{La batterie d'un smartphone}
La batterie d'un smartphone porte les indications : \qty{3,85}{V} ; \qty{4000}{mAh}. On la recharge avec un chargeur de \qty{18}{W} ; le rendement
de la charge est de \qty{85}{\percent}.
\begin{donnees}
$e=\qty{1,60e-19}{C}$ ; \qty{1}{mAh} est la charge transportée par un courant de \qty{1}{mA} pendant une heure.
\end{donnees}
\begin{questions}
\item Que représente l'indication \qty{4000}{mAh} ? La convertir en coulombs.
\item Calculer le nombre d'électrons correspondant.
\item Calculer l'énergie stockée dans la batterie, en joules puis en wattheures.
\item Estimer la durée d'une recharge complète.
\end{questions}
\end{exercice}
\end{blocLivre}
\begin{blocLivre}{ex-14-8}
\begin{exercice}[2]{Ne jamais court-circuiter une pile}
Une pile plate a pour caractéristiques $E=\qty{9,0}{V}$ et $r=\qty{2,0}{\ohm}$.
\begin{questions}
\item Calculer l'intensité de court-circuit et la puissance alors dissipée dans la pile. Conclure.
\item On branche la pile aux bornes d'un conducteur ohmique de \qty{16}{\ohm}. Montrer que l'intensité vaut \qty{0,50}{A}.
\item Calculer la tension aux bornes de la pile et la puissance reçue par le conducteur ohmique.
\item Faire le bilan de puissance de la pile et calculer son rendement.
\end{questions}
\end{exercice}
\end{blocLivre}
\begin{blocLivre}{ex-14-9}
\begin{exercice}[2]{Démarrer par grand froid}
Une batterie de voiture a une tension à vide $E=\qty{12,6}{V}$ et une résistance interne $r=\qty{0,010}{\ohm}$. Au démarrage, le démarreur appelle
un courant de \qty{200}{A}.
\begin{questions}
\item Calculer la tension aux bornes de la batterie pendant le démarrage.
\item Calculer la puissance reçue par le démarreur et celle dissipée dans la batterie.
\item Par grand froid, la résistance interne atteint \qty{0,030}{\ohm}. Calculer la nouvelle tension aux bornes.
\item Le démarreur ne fonctionne plus en dessous d'environ \qty{9}{V}. Expliquer les pannes de démarrage en hiver.
\end{questions}
\end{exercice}
\end{blocLivre}
\begin{blocLivre}{ex-14-10}
\begin{exercice}[2]{Le bilan de puissance d'un circuit}
\schema{14-ex10}
\begin{questions}
\item En écrivant que la tension $U$ est la même pour le générateur et pour la lampe, montrer que $I=\dfrac{E}{r+R}$ et la calculer.
\item Calculer $U$, puis la puissance reçue par la lampe.
\item Calculer la puissance $E\,I$ et la puissance dissipée dans le générateur. Vérifier le bilan de puissance.
\item Calculer le rendement du générateur.
\end{questions}
\end{exercice}
\end{blocLivre}
\begin{blocLivre}{ex-14-11}
\begin{exercice}[2]{Un panneau photovoltaïque}
Un panneau photovoltaïque de \qty{1,7}{m^2} reçoit, en plein soleil, une puissance lumineuse de \qty{1000}{W} par mètre carré. Il délivre alors
une tension de \qty{34}{V} et un courant de \qty{10}{A}.
\begin{questions}
\item Calculer la puissance lumineuse reçue et la puissance électrique fournie.
\item Calculer le rendement du panneau. Représenter sa chaîne de puissance.
\item Sur une année, le panneau fournit autant d'énergie qu'en 1300 heures de plein soleil. Calculer l'énergie produite.
\item Une famille consomme \qty{2500}{kWh} d'électricité par an. Combien de panneaux faudrait-il ?
\end{questions}
\end{exercice}
\end{blocLivre}
\rubrique{Objectif DS}
\begin{blocLivre}{ex-14-12}
\begin{exercice}[3]{Le chauffe-eau électrique}
\begin{docu}{1}{L'appareil}
Un chauffe-eau électrique de \qty{200}{L} contient une résistance chauffante de \qty{3,0}{kW}, alimentée sous \qty{230}{V}. Il fonctionne la nuit, en
heures creuses, où le kilowattheure coûte 0,20~€.
\end{docu}
\begin{docu}{2}{Chauffer de l'eau}
Pour élever de \qty{1}{\degreeCelsius} la température de \qty{1}{kg} d'eau, il faut \qty{4180}{J}. Un litre d'eau a une masse de \qty{1}{kg}. Le rendement
du chauffe-eau, qui perd un peu de chaleur, vaut \qty{90}{\percent}.
\end{docu}
\begin{questions}
\item Calculer l'énergie nécessaire pour chauffer l'eau de \qty{15}{\degreeCelsius} à \qty{65}{\degreeCelsius}, en joules puis en kilowattheures.
\item Calculer l'intensité du courant et la valeur de la résistance chauffante.
\item \initiative Calculer la durée de la chauffe, en tenant compte du rendement.
\item Calculer le coût d'une chauffe complète.
\end{questions}
\end{exercice}
\end{blocLivre}
\begin{blocLivre}{ex-14-13}
\begin{exercice}[3]{La voiture électrique}
\begin{docu}{1}{Fiche technique}
Batterie : \qty{52}{kWh} sous \qty{400}{V}. Puissance maximale du moteur : \qty{100}{kW}. Consommation moyenne : \qty{15}{kWh} pour \qty{100}{km}.
\end{docu}
\begin{docu}{2}{Les recharges}
À domicile, une borne murale délivre \qty{7,4}{kW} avec un rendement de \qty{90}{\percent}. Sur autoroute, une borne rapide de \qty{100}{kW} recharge la
batterie jusqu'à \qty{80}{\percent}.
\end{docu}
\begin{questions}
\item Estimer l'autonomie de la voiture.
\item Calculer l'intensité du courant débité par la batterie lorsque le moteur fournit sa puissance maximale.
\item Calculer la durée d'une recharge complète à domicile.
\item \initiative Estimer la durée d'une recharge rapide de 0 à \qty{80}{\percent}, en supposant la puissance constante.
\end{questions}
\end{exercice}
\end{blocLivre}
\begin{blocLivre}{ex-14-14}
\begin{exercice}[3]{Pourquoi la haute tension ?}
\begin{docu}{1}{Le transport de l'électricité}
Une centrale envoie une puissance de \qty{1000}{MW} vers une ville située à \qty{100}{km}. Les câbles de la ligne ont une résistance de
\qty{0,030}{\ohm} par kilomètre. Le réseau français transporte l'électricité sous \qty{400}{kV}, avant de l'abaisser près des consommateurs.
\end{docu}
\begin{questions}
\item Calculer l'intensité du courant dans la ligne sous \qty{400}{kV}, puis la puissance perdue par effet Joule. Quelle fraction de la puissance
transportée représente-t-elle ?
\item Reprendre le calcul si la ligne fonctionnait sous \qty{20}{kV}. Conclure.
\item \initiative Par quel facteur les pertes sont-elles divisées quand la tension est multipliée par 20 ?
\item Proposer une autre façon de réduire les pertes, et en citer l'inconvénient.
\end{questions}
\end{exercice}
\end{blocLivre}
\begin{blocLivre}{ex-14-15}
\begin{exercice}[3]{Le vrai prix de l'énergie d'une pile}
\begin{docu}{1}{Deux sources pour une lampe torche}
Une pile alcaline AA (\qty{1,5}{V} ; \qty{2,5}{Ah}) coûte 2~€ et ne sert qu'une fois. Un accumulateur rechargeable de même format (\qty{1,2}{V} ;
\qty{2,0}{Ah}) coûte 5~€ et supporte 500 recharges. Le kilowattheure du réseau électrique coûte 0,25~€.
\end{docu}
\begin{questions}
\item Calculer l'énergie stockée dans la pile alcaline, en wattheures puis en joules.
\item Combien de temps peut-elle alimenter une lampe de \qty{0,50}{W} ?
\item \initiative Calculer le prix du kilowattheure fourni par cette pile et le comparer à celui du réseau.
\item Calculer le prix du kilowattheure fourni par l'accumulateur sur toute sa durée de vie (on néglige le coût des recharges). Conclure.
\end{questions}
\end{exercice}
\end{blocLivre}
